\documentclass[10pt]{amsart}
\usepackage[a4paper,margin=23mm]{geometry}
\usepackage{amsmath,amssymb,amsthm,mathtools,hyperref,graphicx,xspace,letltxmacro,xpatch}
\usepackage[english]{babel}
\newtheorem{theo}{Theorem}[section]
\newtheorem{lemm}[theo]{Lemma}\newtheorem{prop}[theo]{Proposition}\newtheorem{coro}[theo]{Corollary}\newtheorem{conj}[theo]{Conjecture}
\theoremstyle{definition}\newtheorem{defi}[theo]{Definition}\newtheorem{rema}[theo]{Remark}\newtheorem{exam}[theo]{Example}\newtheorem{exem}[theo]{Example}\newtheorem{nota}[theo]{Notation}
\newcommand{\ie}{i.e.}\newcommand{\eg}{e.g.}\newcommand{\resp}{resp.}\newcommand{\skpt}{\par\smallskip\noindent}
\emergencystretch=3em\pagestyle{plain}
%\usepackage{enumerate}
\usepackage{dsfont}

\DeclareMathOperator{\modd}{mod}
\newcommand{\bbT}{\mathbb{T}}
\newcommand{\Z}{\mathbb{Z}}
\newcommand*\R{\mathbb{R}}
\newcommand*\N{\mathbb{N}}
\newcommand*\bbL{\mathbb{L}}
\newcommand{\PP}{\mathbb{P}}
\newcommand{\bbE}{\mathbb{E}}
\newcommand{\E}[1]{\bbE\left[#1\right]}

\newcommand*\T{\mathcal{T}}
\newcommand*\I{\mathcal{I}}
\newcommand*\J{\mathcal{J}}
\newcommand*\C{\mathcal{C}}
\newcommand*\A{\mathcal{A}}
\newcommand*\B{\mathcal{B}}
\newcommand*\D{\mathcal{D}}
\newcommand*\F{\mathcal{F}}
\newcommand*\K{\mathcal{K}}
\newcommand*\clH{\mathcal{H}}
\renewcommand*\P{\mathcal{P}}
\newcommand*\G{\mathcal{G}}
\newcommand{\Aca}{\mathcal{A}}
\newcommand{\Fca}{\mathcal{F}}
\newcommand*\clN{\mathcal{N}}


\newcommand{\pr}[1]{\left(#1\right)}
\newcommand{\ind}[1]{{\mathds{1}}_{#1}}

\newcommand{\norm}[1]{\left\lVert #1\right\rVert}
\newcommand{\abs}[1]{\left\lvert #1\right\rvert}
\newcommand{\ens}[1]{\left\{#1\right\}}


%%%%%%%%%%%%%%%%%%%%%%%%%%%%%%%%%%%%%%%%%%%%%%%%%%%%%%%%%%%%%%%%%%%%%%%%%%%%%%%%%%%%%%%%%%%%%%
\renewcommand*{\to}{\mathchoice{\longrightarrow}{\rightarrow}{\rightarrow}{\rightarrow}}

%~ This fixes spacing issues around \left( and \right) and other similar symbols
\let\originalleft\left 
\let\originalright\right
\renewcommand{\left}{\mathopen{}\mathclose\bgroup\originalleft}
\renewcommand{\right}{\aftergroup\egroup\originalright}

%Better widebar

\makeatletter
\let\save@mathaccent\mathaccent
\newcommand*\if@single[3]{%
  \setbox0\hbox{${\mathaccent"0362{#1}}^H$}%
  \setbox2\hbox{${\mathaccent"0362{\kern0pt#1}}^H$}%
  \ifdim\ht0=\ht2 #3\else #2\fi
  }
%The bar will be moved to the right by a half of \macc@kerna, which is computed by amsmath:
\newcommand*\rel@kern[1]{\kern#1\dimexpr\macc@kerna}
%If there's a superscript following the bar, then no negative kern may follow the bar;
%an additional {} makes sure that the superscript is high enough in this case:
\newcommand*\widebar[1]{\@ifnextchar^{{\wide@bar{#1}{0}}}{\wide@bar{#1}{1}}}
%Use a separate algorithm for single symbols:
\newcommand*\wide@bar[2]{\if@single{#1}{\wide@bar@{#1}{#2}{1}}{\wide@bar@{#1}{#2}{2}}}
\newcommand*\wide@bar@[3]{%
  \begingroup
  \def\mathaccent##1##2{%
%Enable nesting of accents:
    \let\mathaccent\save@mathaccent
%If there's more than a single symbol, use the first character instead (see below):
    \if#32 \let\macc@nucleus\first@char \fi
%Determine the italic correction:
    \setbox\z@\hbox{$\macc@style{\macc@nucleus}_{}$}%
    \setbox\tw@\hbox{$\macc@style{\macc@nucleus}{}_{}$}%
    \dimen@\wd\tw@
    \advance\dimen@-\wd\z@
%Now \dimen@ is the italic correction of the symbol.
    \divide\dimen@ 3
    \@tempdima\wd\tw@
    \advance\@tempdima-\scriptspace
%Now \@tempdima is the width of the symbol.
    \divide\@tempdima 10
    \advance\dimen@-\@tempdima
%Now \dimen@ = (italic correction / 3) - (Breite / 10)
    \ifdim\dimen@>\z@ \dimen@0pt\fi
%The bar will be shortened in the case \dimen@<0 !
    \rel@kern{0.6}\kern-\dimen@
    \if#31
      \overline{\rel@kern{-0.6}\kern\dimen@\macc@nucleus\rel@kern{0.4}\kern\dimen@}%
      \advance\dimen@0.4\dimexpr\macc@kerna
%Place the combined final kern (-\dimen@) if it is >0 or if a superscript follows:
      \let\final@kern#2%
      \ifdim\dimen@<\z@ \let\final@kern1\fi
      \if\final@kern1 \kern-\dimen@\fi
    \else
      \overline{\rel@kern{-0.6}\kern\dimen@#1}%
    \fi
  }%
  \macc@depth\@ne
  \let\math@bgroup\@empty \let\math@egroup\macc@set@skewchar
  \mathsurround\z@ \frozen@everymath{\mathgroup\macc@group\relax}%
  \macc@set@skewchar\relax
  \let\mathaccentV\macc@nested@a
%The following initialises \macc@kerna and calls \mathaccent:
  \if#31
    \macc@nested@a\relax111{#1}%
  \else
%If the argument consists of more than one symbol, and if the first token is
%a letter, use that letter for the computations:
    \def\gobble@till@marker##1\endmarker{}%
    \futurelet\first@char\gobble@till@marker#1\endmarker
    \ifcat\noexpand\first@char A\else
      \def\first@char{}%
    \fi
    \macc@nested@a\relax111{\first@char}%
  \fi
  \endgroup
}
\makeatother

\let\oldbar\bar
\renewcommand*{\bar}[1]{{\mathchoice{\widebar{#1}}{\widebar{#1}}{\widebar{#1}}{\oldbar{#1}}}}
%%%%%%%%%%%%%%%%%%%%%%%%%%%%%%%%%%%%%%%%%%%%%%%%%
\graphicspath{{./figures/}}

\newcommand*{\mk}{\mkern -1mu}
\newcommand*{\Mk}{\mkern -2mu}
\newcommand*{\mK}{\mkern 1mu}
\newcommand*{\MK}{\mkern 2mu}

%\hypersetup{urlcolor=purple, linkcolor=blue, citecolor=red}

\newcommand*{\relabel}{\renewcommand{\labelenumi}{(\theenumi)}}
\newcommand*{\romanenumi}{\renewcommand*{\theenumi}{\roman{enumi}}\relabel}
\newcommand*{\Romanenumi}{\renewcommand*{\theenumi}{\Roman{enumi}}\relabel}
\newcommand*{\alphenumi}{\renewcommand*{\theenumi}{\alph{enumi}}\relabel}
\newcommand*{\Alphenumi}{\renewcommand*{\theenumi}{\Alph{enumi}}\relabel}
%%%%%%%%%%%%%%%%%%%%%%%%%%%%%%%%%%%%%%%%%%%%%%%%%%%%%%%%%%%%%%%%%%%%%%%%%%%%%%%%%%%%%%%%%%%%%%%%%%%%%%%%%%%%%%%%%%%%%%%%%%%%%%%%%%%%%%%%%%%%%%%

\newenvironment{enonce*}[2][]{\par\smallskip\noindent\textbf{#2.}\ }{\par\smallskip}
\title{Gaussian-chaos central limit law for product-type three-dimensional martingale-difference fields}
\author{Giraudo, Davide, Lesigne, Emmanuel, Volný, Dalibor}
\date{doi:10.5802/ahl.264; CC BY 4.0}
\begin{document}
\maketitle
\noindent\textit{Editorial note.} Only the selected theorem is counted. Complete original mathematical body and appendix retained as uncounted context. Bibliography text extracted from the retained publisher HTML references; macros preserved with amsart wrapper.\par\medskip
\noindent\textit{Editorial note.} Source-native cite keys rebound to original publisher bibliography entries; mapping retained. Missing cedram proofc/enonce-star rendering environments supplied as applicable.\par\medskip
\noindent\textit{Allowed standard prerequisites (editorial).}\begin{itemize}
\item Standard functional analysis, measure theory and the author-cited classical theorems. Standard prior results stay cited; complete source-local proofs and definitions retained.
\end{itemize}\medskip

\section{Introduction}\label{introduction}

We study here limit theorems of CLT (Central Limit Theorem) type for stationary multiparameters martingale indexed by $\Z^d$. {A part of our results is valid in any dimension $d\geq2$ (Sections~\ref{product-type factor} to~\ref{product-type CLT}), while another part is limited to dimension $d=2$ (Section~\ref{dim2}).}

The framework of our work has been already described in several articles (see \cite{MR3913270} and references therein) and we recall it now.

For $\underline{i}=(i_1,i_2,\dots,i_d)$ and $\underline{j}=(j_1,j_2,\dots,j_d)\in
\Z^d$, we use the two following notations:
\begin{itemize}
\item $\underline{i}\wedge\underline{j}=(\min(i_1,j_1),\min(i_2,j_2),\dots,
\min(i_d,j_d))$;
\item
$\underline{i}\leq\underline{j}\Leftrightarrow i_1\leq j_1$ and $i_2\leq j_2$ and,\dots, and $i_d\leq j_d$.
\end{itemize}
Moreover $\underline{0}$ stands for the element $(0,0,\dots,0)$ of $\Z^d$.


We consider a $\Z^d$ measure preserving action $T=(T_{\underline{i}})_{\underline{i}\in\Z^d}$ on a probability space $(\Omega,\A,\mu)$, equipped with a \emph{completely commuting invariant filtration} $(\F_{\underline{i}})_{\underline{i}\in\Z^d}$, that is a family of sub-$\sigma$-algebras of $\A$ satisfying for all $\underline{i}$ and $\underline{j}$ in $\Z^d$,
\begin{enumerate}
\item \label{propriete_invariante}
$\F_{\underline{i}}=T_{-\underline{i}}\F_{\underline{0}}$;
\item\label{propriete_commutativite}
For all integrable function $f$, $\bbE[\bbE[f\mid\F_{\underline{i}}]\mid\F_{\underline{j}}]=\bbE[f\mid\F_{\underline{i}\wedge\underline{j}}]$.
\end{enumerate}
Note that property~\eqref{propriete_commutativite} implies that $\F_{\underline{i}}\cap\F_{\underline{j}}=\F_{\underline{i}\wedge\underline{j}}$ and in particular $\F_{\underline{i}}\subset\F_{\underline{j}}$ when $\underline{i}\leq \underline{j}$.

We will use classical notations for the limit sub-$\sigma$-algebras when parameters go to $\pm\infty$:
\[
\F_{-\infty,i_2,\dots,i_d}=\bigcap_{i_1\in\Z}\F_{i_1,i_2,\dots,i_d}\quad\text{and}\quad\F_{\infty,i_2,\dots,i_d}=\bigvee_{i_1\in\Z}\F_{i_1,i_2,\dots,i_d},
\]
and so on.

A \emph{field of martingale differences adapted to the filtration} $(\F_{\underline{i}})_{\underline{i}\in\Z^d}$ is a field of random variables $(X_{\underline{i}})_{\underline{i}\in\Z^d}$ of the type
\[
X_{\underline{i}}=f\circ T_{\underline{i}}
\]
where $f\in \bbL^2(\Omega,\F_{\underline{0}},\mu)$ satisfies
\begin{equation}\label{eq:cond-mart}
\E{f\mid\F_{-1,\infty,\dots,\infty}}=\E{f\mid\F_{\infty,-1,\infty,\dots,\infty}}=\dots=\E{f\mid\F_{\infty,\infty,\dots,\infty,-1}}=0.
\end{equation}
(We will say simply that $f$ is a martingale difference adapted to the filtration $(\F_{\underline{i}})$.)

From a previous article~\cite{MR3913270}, we know that, for any such field, we have convergence in law of
\begin{equation}\label{eq:sommes_partielles_normalisees}
\frac1{\sqrt{n_1n_2\dots n_d}}\sum_{i_1=1}^{n_1}\;\sum_{i_2=1}^{n_2}\dots\sum_{i_d=1}^{n_d} X_{i_1,i_2,\dots,i_d}
\end{equation}
when $\min(n_1,n_2,\dots,n_d)$ goes to infinity.

Moreover the limit is normal as soon as one of the $d$ transformations
\[
T_{1,0,\dots,0},
\; T_{0,1,0,\dots,0},\dots \text{ or }\;T_{0,0,\dots,0,1}
\]
is ergodic (this was already established in~\cite{MR3427925}), but some easy examples show that the limit is not normal in general. Indeed, following the example in~\cite{MR3222815}, we can take $X_{i_1,i_2,\dots,i_d}=X^{(1)}_{i_1}X^{(2)}_{i_2},\dots, X^{(d)}_{i_d}$, where $(X^{(1)}_i)_{i\in\Z}$, $(X^{(2)}_i)_{i\in\Z}$,\dots and $(X^{(d)}_i)_{i\in\Z}$ are $d$ mutually independent i.i.d.\ sequences of standard normal random variables, so that for each $(n_1,n_2, \dots,n_d)$, the random variable defined by~\eqref{eq:sommes_partielles_normalisees} has the same distribution as the product of $d$ independent random variables having standard normal distribution. (Note that this example is produced by an ergodic $\Z^d$-action.)


Let us also mention the papers~\cite{MR3504508,MR3869881}, which bring further results and examples.

{In all the sequel, we suppose that the $\Z^d$-action $T$ is ergodic on $(\Omega,\A,\mu)$. Moreover, we suppose that $\F_{\infty,\infty,\dots,\infty}=\A$ which does not cost anything since the whole process we are interested in is $\F_{\infty,\infty,\dots,\infty}$-measurable.}

Here is the general organization of this article.

In Section~\ref{product-type factor} we describe a particular factor $\I$ on which the {$\Z^d$-action} is of \emph{product type}. In Section~\ref{proj-mart}, we show that the martingale structure is preserved by projection on this factor. For a product type action, the possible limit distributions of~\eqref{eq:sommes_partielles_normalisees} are fully understood, as described in Section~\ref{product-type CLT}.

In Section~\ref{dim2}, where we restrict to the case of $\Z^2$-actions, we study what can happen on the orthocomplement of the factor $\I$. In particular we obtain the following results:
\begin{itemize}
\item if the transformation $T_{1,0}$ acting on the factor of $T_{0,1}$-invariants has zero entropy (or if the transformation $T_{0,1}$ acting on the factor of $T_{1,0}$-invariants has zero entropy), then for any square integrable martingale difference the limit distribution in the CLT is normal.

\item if the transformation $T_{1,0}$ acting on the factor of $T_{0,1}$-invariants and the transformation $T_{0,1}$ acting on the factor of $T_{1,0}$-invariants have positive entropies, then there exists a square integrable martingale difference for which the limit distribution in the CLT is not normal.
\end{itemize}
{Proposition~\ref{queue-condition} in Section~\ref{queue-property} gives a sufficient condition for the limit distribution to be normal. In our opinion, the most interesting questions we leave opened in our study concern this proposition: weakening of the moment condition ($\bbL^4$); extension to higher dimension actions.}

\subsection*{Warnings}

\begin{itemize}
\item In the Sections~\ref{product-type factor}, \ref{proj-mart} and~\ref{product-type CLT}, stated results are valid for $\Z^d$ actions for arbitrary $d\geq 2$. We write down statements and proofs in the case $d=3$ because it seems sufficient to perfectly understand the general case and it makes for lighter drafting.
\item On the contrary, Section~\ref{dim2} is specific to the case $d=2$.
\item All identities between $\sigma$-algebras have to be understood modulo the measure.
\end{itemize}
%\end{enonce*}

\subsection*{Acknowledgments}
The authors would like to thank the referee for a careful reading of the manuscript and the suggestions of improvement.

\section{A factor of product type}\label{product-type factor}
{Let us denote by $\I_1$, $\I_2$, $\I_3$ the $\sigma$-algebras of, respectively $T_{1,0,0}$, $T_{0,1,0}$, $T_{0,0,1}$ invariant sets in $\A$.

Let us denote by $\overline{\I_1}$, $\overline{\I_2}$, $\overline{\I_3}$ the $\sigma$-algebras of invariant sets under, respectively, the $\Z^{2}$-actions $(T_{0,i_2,i_3})$, $(T_{i_1,0,i_3})$, $(T_{i_1,i_2,0})$. (For an arbitrary value of $d$, we consider here $\Z^{d-1}$-actions.) In other words, we have
\[
\overline{\I_t}=\bigcap_{1\leq s\leq 3,\,s\neq t}\I_s, \quad 1\leq t\leq 3.
\]
(Note that in the particular case $d=2$, we have $\overline{\I_1}=\I_2$.)

Finally we note $\I$ the $\sigma$-algebra generated by the union of the $\overline{\I_t}\,$:
\[
\I=\bigvee_{t=1}^3\overline{\I_t}\;.
\]
}

Note that all these $\sigma$-algebras are invariant under the action $T$, hence they define factors of the dynamical system $(\Omega,\A,\mu,T)$.

\begin{prop}\label{indep}
The $\sigma$-algebras $\overline{\I_t}$, $1\leq t\leq 3$ are mutually independent.
\end{prop}

As a consequence of this proposition we can state that the action $T$ on the probability space $(\Omega,\I,\mu)$ is of \emph{product type}, which means that it is isomorphic to a~$\Z^3$-action defined on the product of $3$ probability spaces $(\Omega_t,\A_t,\mu_t)$ by a formula of the type
\[
T_{i_1,i_2,i_3}(\omega_1,\omega_2,\omega_3)=
\left(T_1^{i_1}\omega_1,T_2^{i_2}\omega_2,T_3^{i_3}\omega_3\right)
\]
where, for each $t$, $T_t$ is an invertible ergodic measure preserving transformation of $(\Omega_t,\A_t,\mu_t)$.\footnote{The fact that the sub-$\sigma$-algebras $\overline{\I_1}$, $\overline{\I_2}$ et $\overline{\I_3}$ are independent and generate the $\sigma$-algebra $\I$ tells us that the measure algebra dynamical systems
\[
(\I,\mu,T) \quad\text{and}\quad (\overline{\I_1},\mu,T) \times (\overline{\I_2},\mu,T) \times (\overline{\I_3},\mu,T)
\]
are isomorphic (in the sense of~\cite[Definition~2.14]{MR1958753} in Glasner book, which corresponds to the notion of conjugacy in the classical Billingsley book~\cite{MR192027}). The methods presented in~\cite[Chapter~2]{MR1958753}, can be used to build Lebesgue measure preserving dynamical systems models for these measure algebra dynamical systems.}


Proposition~\ref{indep} is a straightforward consequence of the following general lemma, brought to our attention by the referee of this article.

\begin{lemm}\label{lem-indep}
Let $G$ and $H$ be two measure preserving group actions on a~probability space $(X,\A,\mu)$ which commute and such that the action generated by both $G$ and $H$ is ergodic. Then the $\sigma$-algebras of, respectively, $G$-invariant sets and\linebreak $H$-invariant sets are independent.
\end{lemm}

\begin{proof}[Proof of Lemma~\ref{lem-indep}]
Denote by $\G$ and $\clH$, respectively, the $\sigma$-algebras of \linebreak $G$-invariant sets and $H$-invariant sets. Let $u$ be a bounded $\G$-measurable function on~$X$. Since the actions commute the $\sigma$-algebra $\clH$ is $G$-invariant. Since $u$ and $\clH$ are both $G$-invariant, $\bbE[u\mid\clH]$ is also $G$-invariant. The ergodicity of the action generated by $G$ and $H$ means that $\G\cap\clH$ is trivial, and we conclude that $\bbE[u\mid\clH]$ is constant. This establishes the independence property.
\end{proof}


\section{Martingale property preserved by projection}\label{proj-mart}

Our aim here is to show that the martingale property is preserved by projection on the factor~$\I$. We begin by a general abstract result, which is stated for a $\Z^{2}$-measure preserving action. (For an arbitrary value of $d$, we have to consider here a~$\Z^{d-1}$-action.) 

Let $(X,\B,\nu)$ be a probability space and $U=(U_{\underline{i}})$ a $\Z^2$-measure preserving action on this space. We denote by $\J$ the $\sigma$-algebra of $U$-invariant elements of the $\sigma$-algebra~$\B$.

We denote by $\F$ a $U$-invariant sub-$\sigma$-algebra of $\B$, meaning that $U_{\underline{i}}(\F)=\F$ for all $\underline{i}\in\Z^2$; we denote by $\C$ a sub-$\sigma$-algebra of $\F$.

\begin{prop}\label{pro}
\leavevmode
\begin{enumerate}\romanenumi
\item\label{assertion_EfF_mea}
For all $f\in \bbL^2(\J\vee\C)$, $\bbE[f\mid\F]$ is $\J\vee\C$-measurable.
\item\label{assertion_F_mes_J_vee_C}
For all $f\in \bbL^2(\F)$, $\bbE[f\mid\J\vee\C]$ is $\F$-measurable.
\end{enumerate}
\end{prop}

\begin{coro}\label{co}
The conditional expectations with respect to $\F$ and to $\J\vee\C$ are commuting: for all $f\in \bbL^2(\B)$,
\[
\bbE\bigl[\E{f\mid\J\vee\C}\,\bigm|\,\F\bigr]=\bbE\bigl[\E{f\mid\F}\,\bigm|\,\J\vee\C\bigr]=\E{f\,\middle|\,\F\cap(\J\vee\C)}.
\]
\end{coro}


\begin{proof}[Proof of Corollary~\ref{co}]
Suppose that $\D$ and $\D'$ are two sub-$\sigma$-algebras such that, for all $f\in \bbL^2(\D)$, $\bbE[f\mid\D']$ is $\D$-measurable. By ordinary properties of projections, we have, for any $f\in \bbL^2(\B)$
\begin{align*}
\E{f\mid\D\cap\D'} &=\E{\E{\E{f\mid\D}\mid\D'}\mid\D\cap\D'}
\\
\intertext{which implies}
\E{f\mid\D\cap\D'} &=\E{\E{f\mid\D}\mid\D'}
\end{align*}
since $\bbE[\bbE[f\mid\D]\mid\D']$ is $\D\cap\D'$-measurable.
\end{proof}


\begin{proof}[Proof of Proposition~\ref{pro}]
Let $f\in \bbL^2(\F)$; by the ergodic theorem, we have
\[
\E{f\mid\J}=\lim_{\ell,m\to\infty}\frac1{\ell m}\sum_{i=1}^\ell \sum_{j=1}^mf\circ U_{i,j},
\]
hence $\bbE[f\mid\J]$ is $\F$-measurable, since $\F$ is $U$-invariant.

For any $f\in \bbL^2(\B)$, we have
\[
\E{f\mid\F}\circ U_{i,j}=\E{f\circ U_{i,j}\,\middle|\, U_{i,j}^{-1}(\F)}=\E{f\circ U_{i,j}\,\middle|\,\F},
\]
hence if $f$ is $\J$-measurable, then $\bbE[f\mid\F]$ is $\J$-measurable.

Consider now $g\in \bbL^2(\J)$ and $h\in \bbL^2(\C)$. Since $\C\subset \F$, we have $
\bbE[gh\mid\F]=\bbE[g\mid\F] h $, so $\bbE[gh\mid\F]$ is $\J\vee\C$-measurable. But the functions of the form $gh$ with $g\in \bbL^\infty(\J)$ and $h\in \bbL^\infty(\C)$ generate a dense subspace of $ \bbL^2(\J\vee\C)$, so that assertion~\eqref{assertion_EfF_mea} is proved.

Consider now $f\in \bbL^2(\F)$, $g\in \bbL^2(\J)$ and $h\in \bbL^2(\C)$. We have
\begin{align*}
\left\langle \E{f\mid\J\vee\C},gh\right\rangle &=\langle f,gh\rangle\quad\text{(because $gh$ is $\J\vee\C$-measurable)},\\
&=\langle f,\E{gh\mid\F}\rangle\quad\text{(because $f$ is $\F$-measurable)},\\
&=\langle f,h\E{g\mid\F}\rangle\quad\text{(because $h$ is $\F$-measurable)},\\
&=\langle \E{f\mid\J\vee\C},h\E{g\mid\F}\rangle \quad\text{(because $\E{g\mid\F}$ is $\J$-measurable)},\\
&=\langle \E{f\mid\J\vee\C},\E{gh\mid\F}\rangle.
\end{align*}
By the density argument, we conclude that, for all $k\in \bbL^2(\J\vee\C)$,
\[
\langle \E{f\mid\J\vee\C},k)\rangle=\langle \E{f\mid\J\vee\C},\;\E{k\mid\F}\rangle.
\]
This identity applied to the function $k= \bbE[f\mid\J\vee\C]$ shows that this function is $\F$-measurable. This is assertion~\eqref{assertion_F_mes_J_vee_C}.
\end{proof}

Another result we need is the following classical lemma (cf for example~\cite{MR262464}) and we give a short proof for the sake of completeness.
\begin{lemm} \label{class}
Let $S$ be a measure preserving transformation of the probability space $(X,\B,\nu)$, and $\F_n = S^{-n}\F_0$ be an increasing filtration in $\B$. Denote by $\K$ the sub-$\sigma$-algebra of $S$ invariant sets. Then $\K \cap \F_{\infty}=\K \cap \F_{-\infty}$.
In particular, if $\K \subset \F_{\infty}$ then $\K \subset \F_{-\infty}$.
\end{lemm}
\begin{proof}[Proof of Lemma~\ref{class}]
Let $f\in \bbL^2(\nu)$. By purely Hilbert space arguments, we know that, in the space $ \bbL^2(X)$,
\[
\lim_{n\to-\infty}\E{f\mid\F_n}=\E{f\mid\F_{-\infty}}
\quad\text{and}\quad \lim_{n\to\infty}\E{f\mid\F_n}=\E{f\mid\F_{\infty}}.
\]
Now suppose that $f$ is invariant under $S$, that is $f$ is $\K$-measurable. Then
\[
\E{f\mid\F_n}=\E{f\circ S^n\mid S^{-n}\F_0}=\E{f\mid\F_0 }\circ S^n.
\]
Thus $\|\bbE[f\mid\F_{-m}]\|_2=\|\bbE[f\mid\F_n]\|_2$ and with $n,m\to\infty$, we obtain $\|\bbE[f\mid\F_{-\infty}]\|_2=\|\bbE[f\mid\F_\infty]\|_2$. Since $\F_{-\infty}\subset\F_\infty$, this implies $
\bbE[f\mid\F_{-\infty}]=\bbE[f\mid\F_{\infty}] $. In particular, if $f$ is $\F_{\infty}$-measurable, then it is $\F_{-\infty}$-measurable, which is what we had to prove.
\end{proof}

We now come back to the situation described in the preceding section where the factor $\I$ is defined.

\begin{theo}\label{proj_mart-diff}
Let $f$ be a martingale difference adapted to the filtration $(\F_{\underline{i}})_{\underline{i}\in\Z^3}$. Then its projection $\bbE[f\mid\I]$ is a martingale difference adapted to the filtration {$(\F_{\underline{i}})_{\underline{i}\in\Z^3}$}, as well adapted to the filtration {$(\I\cap\F_{\underline{i}})_{\underline{i}\in\Z^3}$}. Moreover $f-\bbE[f\mid\I]$ is also a martingale difference adapted to the filtration $(\F_{\underline{i}})_{\underline{i}\in\Z^3}$.
\end{theo}
\begin{proof}
We want to apply Proposition~\ref{pro} to the $\Z^2$-action $U_{j,k}=T_{0,j,k}$ on the space $(\Omega,\A,\mu)$. So we have $\J=\bar \I_1$. We consider also $\C=\overline{\I_2}\vee\overline{\I_3}$ and $\F=\F_{i,\infty,\infty}$, for a given $i$. We know that $\F$ is $U$-invariant and, thanks to Lemma~\ref{class}, we have $\I_1\subset\F$, hence $\C\subset\F$ since $\C\subset\I_1$.

Note that $\J\vee\C=\bar \I_1\vee\overline{\I_2}\vee\overline{\I_3}=\I$.

Now Corollary~\ref{co} tells us that the conditional expectation with respect to $\I$ commutes with the conditional expectations with respect to $\F_{i,\infty,\infty}$. Of course, we can exchange the roles of $i$, $j$ and $k$ and we find as well that the conditional expectation with respect to $\I$ commutes with the conditional expectations with respect to $\F_{\infty,j,\infty}$ and with the conditional expectations with respect to $\F_{\infty,\infty,k}$.

By the complete commuting property of the filtration, we have
\[
\E{\cdot\,\middle|\,\F_{i,j,k}}=\E{\E{\E{\cdot\,\middle|\,\F_{i,\infty,\infty}}\,\middle|\,\F_{\infty,j,\infty}}\,\middle|\,\F_{\infty,\infty,k}},
\]
and we conclude that the conditional expectation with respect to $\I$ commutes with the conditional expectations with respect to $\F_{i,j,k}$: for all $f\in \bbL^2(\A)$, for all $i,j,k\in\Z\cup\{\infty\}$,
\begin{equation}\label{eq:commut}
\E{\E{f\,\middle|\,\F_{i,j,k}}\mid\I}=\E{\E{f\,\middle|\,\I}\,\middle|\,\F_{i,j,k}}=\E{f\,\middle|\,\F_{i,j,k}\cap\I}.
\end{equation}

The first thing we want to see now is that $(\I\cap\F_{i,j,k})$ is a completely commuting invariant filtration. The first point is
\[
T_{-i,-j,-k}(\I\cap\F_{0,0,0})=\I\cap\F_{i,j,k}
\]
which is true thanks to~\eqref{assertion_EfF_mea} and the fact that $\I$ is invariant.
The second point is that, for all integrable function $f$,
\[
\E{\E{f\,\middle|\,\I\cap\F_{i,j,k}}\,\middle|\,\I\cap\F_{i',j',k'}}=\E{f\,\middle|\,\I\cap\F_{\min(i,i'),\min(j,j'),\min(k,k')}},
\]
which, thanks to~\eqref{eq:commut}, can be written
\[
\E{\E{\E{f\mid\F_{i,j,k}}\,\middle|\,\F_{i',j',k'}}\,\middle|\,\I}=\E{\E{f\,\middle|\,\F_{\min(i,i'),\min(j,j'),\min(k,k')}}\,\middle|\,\I}
\]
and is true thanks to~\eqref{assertion_F_mes_J_vee_C}.

The second thing to verify is that if $f$ satisfies the martingale condition~\eqref{eq:cond-mart} then $\bbE[f\mid\I]$ satisfies it. The facts that conditional expectations with respect to $\I$ and $\F_{0,0,0}$ commute and that $f$ is $\F_{0,0,0}$-measurable imply that $\bbE[f\mid\I]$ is $\I\cap\F_{0,0,0}$-measurable. Moreover we have, thanks to~\eqref{eq:commut} and~\eqref{eq:cond-mart}
\[
\E{\E{f\mid\I}\mid\F_{-1,\infty,\infty}}=\E{\E{f\mid\F_{-1,\infty,\infty}}\mid\I}=0
\]
and similarly
\[
\E{\E{f\mid\I}\mid\I\cap\F_{-1,\infty,\infty}}=0.
\]

We conclude that $\bbE[f\mid\I]$ is a martingale difference (for any of the two filtrations), and $f-\bbE[f\mid\I]$ is also a martingale difference adapted to the filtration $(\F_{i,j,k})$, just as difference of two martingale differences.
\end{proof}


\section{Limit law in the case of a product type action}\label{product-type CLT}

In this section, we describe the limit distribution of~\eqref{eq:sommes_partielles_normalisees} in the particular case of a product type action equipped with a \emph{product type filtration}, that is a filtration of the type~\eqref{filtration-produit}, as defined below. We explain in Section~\ref{appli_to_I} that this result applies to our factor $\I$.

\subsection{Convergence theorem}

In this section, we deal with the case when {the probability space $(\Omega,\Aca,\mu)$ has a product structure: $\Omega=\Omega_1\times\Omega_2\times\Omega_3$, $\Aca=\Aca_1\otimes\Aca_2\otimes\Aca_3$ and $\mu=\mu_1\otimes\mu_2\otimes\mu_3$, where each $(\Omega_t,\Aca_t,\mu_t)$ is a probability space, equipped with an ergodic invertible measure preserving transformation $T_t$. We consider the action $T$ of $\Z^3$ on $\Omega$ given by
\begin{equation*}
T_{i_1,i_2,i_3}\pr{\omega_1,\omega_2,\omega_3}=
\pr{T_1^{i_1}\omega_1,T_2^{i_2}\omega_2,T_3^{i_3}\omega_3},\quad i_1,i_2,i_3\in\Z.
\end{equation*}
This type of $\Z^3$ measure preserving action will be called a \emph{product type action}.}

{For each $t\in\{1,2,3\}$, consider a sub-$\sigma$-algebra $\Fca_0^{(t)}$ of $\Aca_t$ such that $T_t\Fca_0^{\pr{t}}\subset \Fca_0^{\pr{t}}$. For each $\underline{i}=(i_1,i_2,i_3)\in\Z^3$, define the sub-$\sigma$-algebra $\Fca_{\underline{i}}$ of $\Aca$ by
\begin{equation}\label{filtration-produit}
\Fca_{\underline{i}}=\pr{T_1^{-i_1}\Fca_0^{\pr{1}}}\otimes
\pr{T_2^{-i_2}\Fca_0^{\pr{2}}}\otimes \pr{T_3^{-i_3}\Fca_0^{\pr{3}}}.
\end{equation}
}

\begin{prop}
The filtration {$(\Fca_{\underline{i}})_{\underline{i}\in\Z^3}$} is a completely commuting invariant filtration.
\end{prop}

We omit the proof of this proposition which is relatively straightforward and similar to the proof of~\cite[Proposition~1.2]{MR4754186}.

In order to investigate the convergence of the partial sum process given by~\eqref{eq:sommes_partielles_normalisees}, we will need to decompose the considered function $f$ as a sum of functions which can be expressed as a product of functions of a single $\omega_t$.
\begin{lemm}\label{lem:base_Hilbert}
Denote by $\Delta$ the set of square integrable functions satisfying the martingale property defined as in~\eqref{eq:cond-mart}, {that is $\Delta$ is the space of functions $f\in \bbL^2(\Fca_{\underline0})$ such that
\begin{equation}\label{multi-ortho}
\E{f\mid\Fca_{-1,\infty,\infty}}=\E{f\,\middle|\,\Fca_{\infty,-1,\infty}}=\E{f\,\middle|\,\Fca_{\infty,\infty,-1}}=0.
\end{equation}
There exist random variables $v_{a_t,t}\in \bbL^2(\A_t)$, $1\leq t\leq 3$, $a_t\in\N$, such that the collection of random variables $(\omega_1,\omega_2,\omega_3)\mapsto v_{a_1,1}\pr{\omega_1}v_{a_2,2}\pr{\omega_2} v_{a_3,3}\pr{\omega_3}$ is a Hilbert basis of the space $\Delta$.}
\end{lemm}
\begin{proof}[Proof of Lemma~\ref{lem:base_Hilbert}]
Denote $H^{(t)}=\bbL^2(\Fca_{0}^{\pr{t}})\ominus
\bbL^2(\Fca_{-1}^{\pr{t}})$.
The three orthogonal decompositions
\[
\bbL^2\pr{\Fca_{0}^{\pr{t}}}=H^{(t)} \oplus \bbL^2\left(\Fca_{-1}^{\pr{t}}\right), \quad 1\leq t\leq 3,
\]
give a decomposition of their tensor product:
\begin{align*}
\bbL^2\left(\Fca_{\underline0}\right)&=\left[\bbL^2\pr{\Fca_{-1}^{\pr{1}}} \otimes \bbL^2\pr{\Fca_{0}^{\pr{2}}}\otimes \bbL^2\pr{\Fca_{0}^{\pr{3}}}\right]\\
&\quad\oplus\left[H^{(1)}\otimes \bbL^2\pr{\Fca_{-1}^{\pr{2}}}\otimes\bbL^2\pr{\Fca_{0}^{\pr{3}}}\right]\\
&\quad\oplus \left[H^{(1)}\otimes H^{(2)}\otimes \bbL^2\pr{\Fca_{-1}^{\pr{3}}}\right]\\
&\quad\oplus\left[H^{(1)}\otimes H^{(2)}\otimes H^{(3)}\right].
\end{align*}
Condition~\eqref{multi-ortho} means that the function $f$ is orthogonal to the third first spaces appearing in this decomposition. In other words, we have
\[
\Delta=H^{(1)}\otimes H^{(2)}\otimes H^{(3)}
\]
and we can choose in each space $H^{(t)}$ a Hilbert base $(v_{a,t})_{a\in\N}$.
\end{proof}


\begin{theo}\label{thm:TLC_product}
We consider a product type $\Z^3$-action as defined at the beginning of the section, equipped with a filtration defined as in \eqref{filtration-produit}. Let $f$ be a square integrable function on $\Omega$ such that $(f\circ T_{\underline{i}})_{\underline{i}\in\Z^3}$ is a martingale difference random field. 

There exists a family of real numbers $(\lambda_{a_1,a_2,a_3}\pr{f})_{a_t\geq 1, 1\leq t\leq 3}$
such that
\[
\sum_{a_t\geq 1, 1\leq t\leq 3}\lambda_{a_1,a_2,a_3}^2\pr{f}<\infty
\]
and such that if $(N_a^{\pr{t}})_{a\geq 1}$, $1\leq t\leq 3$, are $3$
i.i.d.\ and mutually independent sequences of standard normal random variables, then
\begin{equation*}
\frac1{\sqrt{n_1n_2 n_3}}\sum_{i_1=1}^{n_1}\sum_{i_2=1}^{n_2}\sum_{i_3=1}^{n_3} f\circ T_{(i_1,i_2,i_3)}\to \sum_{a_t\geq 1, 1\leq t\leq 3} \lambda_{a_1,a_2,a_3}\pr{f} N_{a_1}^{\pr{1}}N_{a_2}^{\pr{2}} N_{a_3}^{\pr{3}}
\end{equation*}
in distribution as $\min\{n_1,n_2,n_3\}\to\infty$.
\end{theo}

\begin{proof}
We know by Lemma~\ref{lem:base_Hilbert} that we can express $f$ as
\[
f(\omega_1,\omega_2,\omega_3)=\sum_{a,b,c=1}^\infty
\lambda_{a,b,c}\pr{f}
v_{a,1}\pr{\omega_1}v_{b,2}\pr{\omega_2}v_{c,3}\pr{\omega_3},
\]
where the convergence takes place in $\bbL^2\pr{\mu}$.

Define
\begin{equation*}
f_K\pr{\omega_1,\omega_2,\omega_3}:=\sum_{a,b,c=1}^K
\lambda_{a,b,c}\pr{f} v_{a,1}\pr{\omega_1}v_{b,2}\pr{\omega_2}v_{c,3}\pr{\omega_3}.
\end{equation*}
Note that $(f_K\circ T^{i,j,k})_{i,j,k\in\Z}$ is also a martingale difference random field. Suppose that we proved for each $K\geq 1$ Theorem~\ref{thm:TLC_product} with $f$ replaced by $f_K$. By orthogonality of increments, for all $\ell, m, n>0$,
\[
\PP\pr{\abs{\frac1{\sqrt{\ell mn}}\sum_{i=1}^\ell\sum_{j=1}^m\sum_{k=1}^n f\circ T_{i,j,k}-\frac1{\sqrt{\ell mn}}\sum_{i=1}^\ell\sum_{j=1}^m\sum_{k=1}^n f_K\circ T_{i,j,k}}>\varepsilon}
\leq\frac 1{\varepsilon^2}\norm{f-f_K}_2^2
\]
hence
\[
\lim_{K\to\infty}\sup_{\ell,m,n}
\PP\pr{\abs{\frac1{\sqrt{\ell mn}}\sum_{i=1}^\ell\sum_{j=1}^m\sum_{k=1}^n f\circ T_{i,j,k}-\frac1{\sqrt{\ell mn}}\sum_{i=1}^\ell\sum_{j=1}^m\sum_{k=1}^n f_K\circ T_{i,j,k}}>\varepsilon}=0.
\]
Moreover,
\begin{multline*}
\lim_{K\to\infty}\norm{\sum_{a,b,c=1}^\infty \lambda_{a,b,c}\pr{f} N_a^{\pr{1}}N_b^{\pr{2}}N_c^{\pr{3}}-\sum_{a,b,c=1}^K
\lambda_{a,b,c}\pr{f} N_a^{\pr{1}}N_b^{\pr{2}}N_c^{\pr{3}}}_2^2\\ =
\lim_{K\to\infty} \sum_{a,b,c=1}^\infty
\mathds{1}_{\max\ens{a,b,c}\geq K+1 }\;
\lambda_{a,b,c}^2\pr{f}=0
\end{multline*}
hence we would get the conclusion of Theorem~\ref{thm:TLC_product} by an application of~\cite[Theorem~4.2]{MR233396}. We thus have to prove that for each $K$,
\begin{equation}\label{eq:tlc_produit_somme_de_K_elements}
\frac1{\sqrt{\ell mn}}\sum_{i=1}^\ell\sum_{j=1}^m\sum_{k=1}^n f_K\circ T_{i,j,k}\to \sum_{a,b,c=1}^K\lambda_{a,b,c}\pr{f} N_a^{\pr{1}}N_b^{\pr{2}}N_c^{\pr{3}}
\end{equation}
in distribution as $\min\{\ell,m,n\}\to\infty$. By definition of $f_K$,
\begin{multline*}
\frac1{\sqrt{\ell mn}}\sum_{i=1}^\ell\sum_{j=1}^m\sum_{k=1}^n f_K\circ T_{i,j,k}\\
=\sum_{a,b,c=1}^K\lambda_{a,b,c}\pr{f}
\pr{\frac 1{\sqrt{\ell}}\sum_{i=1}^\ell v_{a,1}\circ T_1^i}
\pr{\frac 1{\sqrt{m}}\sum_{j=1}^m v_{b,2}\circ T_2^j}
\pr{\frac 1{\sqrt{n}}\sum_{k=1}^n v_{c,3}\circ T_3^k}.
\end{multline*}
Consider the random vector $V_{\ell,m,n}$ of dimension~$3K$, where the first $K$ entries are $\frac 1{\sqrt{\ell}}\sum_{i=1}^\ell v_{a,1}\circ T_1^i$, $1\leq a\leq K$, the entries of index between $K+1$ and $2K$ are $\frac 1{\sqrt{m}}\sum_{j=1}^m v_{b,2}\circ T_2^j$ and the last $K$ are $\frac 1{\sqrt{n}}\sum_{k=1}^n v_{c,3}\circ T_3^k$. By the Cramer--Wold device, the Billingsley--Ibragimov \emph{Central Limit Theorem} for martingale differences and the fact that $\|v_{a,1}\|_2=\|v_{b,2}\|_2=\|v_{c,3}\|_2=1$, the vector $V_{\ell,m,n}$ converges in distribution as
$\min\{\ell,m,n\}\to\infty$ to
\[
V:=\pr{N_1^{\pr{1}},\dots,N_K^{\pr{1}},
N_1^{\pr{2}},\dots,N_K^{\pr{2}},N_1^{\pr{3}},\dots,N_K^{\pr{3}}},
\]
where $N_a^{\pr{1}}$, $N_b^{\pr{2}}$ and $N_c^{\pr{3}}$ are like in the statement of the theorem. Now, \eqref{eq:tlc_produit_somme_de_K_elements} follows from the continuous mapping theorem, that is, $g(V_{\ell,m,n})\to g\pr{V}$, where $g\colon \R^{3K}\to \R$ is defined as
\begin{equation*}
g\pr{x_1,\dots,x_K,y_1,\dots,y_K,z_1,\dots,z_K} =\sum_{a,b,c=1}^K\lambda_{a,b,c}\pr{f} x_ay_bz_c.
\end{equation*}
This ends the proof of Theorem~\ref{thm:TLC_product}.
\end{proof}

\subsection{Application to our situation} \label{appli_to_I}
We come back to the general situation described in the introduction and we consider the factor $\I$ defined in Section~\ref{product-type factor}. We claim that Theorem~\ref{thm:TLC_product} applies to the martingale difference random field $(\bbE[f\mid\I]\circ T_{\underline{i}})_{\underline{i}\in\Z^3}$. We know already that the action of $\Z^3$ on this factor is of product type (cf. Section~\ref{product-type factor}), and that $\bbE[f\mid\I]$ is a martingale difference adapted to the completely commuting invariant filtration $(\I\cap\F_{\underline{i}})_{\underline{i}\in\Z^3}$ (cf. Theorem~\ref{proj_mart-diff}). We just have to verify that this filtration has the product structure described in formula~\eqref{filtration-produit}. Next proposition will apply to the original $\Z^3$-action restricted to the space $(\Omega,\I,\mu)$ ; the key argument will be that $\I\subset\F_{\infty,\infty,\infty}$.

Consider a product type $\Z^3$ measure preserving action as defined at the beginning of Section~\ref{product-type CLT}. We denote by $\I_1$ (resp. $\I_2$, resp. $\I_3$) the sigma-algebra of $T_{1,0,0}$(resp. $T_{0,1,0}$, resp. $T_{0,0,1}$)-invariant events, that is
\[
\I_1=\{\Omega_1,\emptyset\}\otimes\A_2\otimes\A_{3},\quad \I_2=\A_1\otimes\{\Omega_2,\emptyset\}\otimes\A_{3},\quad \I_1=\A_1\otimes\A_2\otimes\{\Omega_3,\emptyset\}.
\]
We denote by $\overline{\I_1}$ (resp. $\overline{\I_2}$, resp. $\overline{\I_3}$) the sigma-algebra of events invariant under the $\Z^{2}$-action $T_{0,j,k}$ (resp. $T_{i,0,k}$, resp. $T_{i,j,0}$) that is
\[
\overline{\I_t}=\bigcap_{s\neq t}\I_s.
\]
These $\sigma$-algebras $\overline{\I_t}$ are mutually independent and $\A=\bigvee_{t=1}^d\overline{\I_t}$.

\begin{prop}\label{filt-prod}
Suppose that $(\G_{\underline{i}})_{\underline{i}\in\Z^3}$ is a completely commuting invariant filtration such that
\[
\A=\G_{\infty,\infty,\infty}.
\]
Then, for each $\underline{i}=(i_1,i_2,i_3)\in\Z^3$,
\[
\G_{\underline{i}}=\bigvee_{t=1}^3 \left(\G_{\underline{i}}\cap\overline{\I_t}\right).
\]
\end{prop}

Note that the conclusion of the proposition can be written as the identity
\[
\G_{\underline{i}}=\bigvee_{t=1}^3 T_{(t)}^{-i_t}\left(\G_{\underline0}\cap\overline{\I_t}\right),
\]
which, thanks to the product structure, is of the type~\eqref{filtration-produit}.

\begin{lemm} \label{cond-ind}
For all $\underline{i}\in\Z^3$ the sigma-algebras $\overline{\I_t}$ ($1\leq t\leq 3$) are conditionally independent with respect to $\G_{\underline{i}}$.
\end{lemm}

\begin{proof}[Proof of Lemma~\ref{cond-ind}]
By hypothesis, $\I_2\subset\G_{\infty,\infty,\infty}$; thanks to Lemma~\ref{class}, we know that, for all $j\in\Z$, $\I_2\subset\G_{\infty,j,\infty}$. Similarly, for all $k\in\Z$, $\I_3\subset\G_{\infty,\infty,k}$. Thus
\[
\overline{\I_1}=\I_2\cap\I_3\subset\G_{\infty,j,\infty}\cap\G_{\infty,\infty,k}=\G_{\infty,j,k}.
\]
Similarly we have $\overline{\I_2}\subset\G_{i,\infty,k}$ and $\overline{\I_3}\subset\G_{i,j,\infty}$, for all $i,j,k\in\Z$.

Since the $\sigma$-algebra $\G_{i,\infty,\infty}$ is $T_{0,j,k}$-invariant we have:
\begin{equation}\label{arg1}
\text{if }u\in \bbL^2\left(\overline{\I_1}\right) \text{ then } \E{u\,\middle|\,\G_{i,\infty,\infty}}\in \bbL^2\left(\overline{\I_1}\right).
\end{equation}

Let $u\in \bbL^2(\overline{\I_1})$, $v\in \bbL^2(\overline{\I_2})$ and $w\in \bbL^2(\overline{\I_3})$. Using the complete commuting property of the filtration $\G$, we write:
\[
\E{uvw\,\middle|\,\G_{i,j,k}}=\bbE\left[\E{uvw\,\middle|\,\G_{i,\infty,\infty}}\,\middle|\,\G_{\infty,j,k}\right].
\]
As we saw just above, $v$ and $w$ are $\G_{i,\infty,\infty}$-measurable. Hence
\[
\E{uvw\,\middle|\,\G_{i,j,k}}=\bbE\left[vw\E{u\,\middle|\,\G_{i,\infty,\infty}}\,\middle|\,\G_{\infty,j,k}\right],
\]
But we know by~\eqref{arg1} that $\bbE[u\mid\G_{i,\infty,\infty}]$ is $\overline{\I_1}$-measurable, hence $\G_{\infty,j,k}$-measurable, and we obtain
\begin{equation*}\label{arg2}
\E{uvw\,\middle|\,\G_{i,j,k}}=\E{u\,\middle|\,\G_{i,\infty,\infty}}\;\E{vw\,\middle|\,\G_{\infty,j,k}}.
\end{equation*}

Repeating similar arguments, we can write
\begin{align*}
\E{vw\,\middle|\,\G_{\infty,j,k}}
&=\bbE\left[\E{vw\,\middle|\,\G_{\infty,j,\infty}}\,\middle|\,\G_{\infty,\infty,k}\right]\\
&=\bbE\left[w\E{v\,\middle|\,\G_{\infty,j,\infty}}\,\middle|\,\G_{\infty,\infty,k}\right]\\
&=\E{v\,\middle|\,\G_{\infty,j,\infty}}\;\E{w\,\middle|\,\G_{\infty,\infty,k}}.
\end{align*}
We obtain
\[
\E{uvw\,\middle|\,\G_{i,j,k}}=\E{u\,\middle|\,\G_{i,\infty,\infty}}\;\E{v\,\middle|\,\G_{\infty,j,\infty}}\;\E{w\,\middle|\,\G_{\infty,\infty,k}}.
\]
In particular
\[
\E{u\,\middle|\,\G_{i,j,k}}=\E{u\,\middle|\,\G_{i,\infty,\infty}},\
\E{v\,\middle|\,\G_{i,j,k}}=\E{v\,\middle|\,\G_{\infty,j,\infty}}
\]
and
\[
\E{w\,\middle|\,\G_{i,j,k}}=\E{w\,\middle|\,\G_{\infty,\infty,k}}
\]
and we conclude that
\[
\E{uvw\,\middle|\,\G_{i,j,k}}=\E{u\,\middle|\,\G_{i,j,k}}\;\E{v\,\middle|\,\G_{i,j,k}}\E{w\,\middle|\,\G_{i,j,k}}.\qedhere
\]
\end{proof}
\begin{proof}[Proof of Proposition~\ref{filt-prod}]
Fix $(i,j,k)\in\Z^3$. The conclusion of the proof of Lemma~\ref{cond-ind} shows that, for any function of the type $f=uvw$ with $u\in \bbL^2(\overline{\I_1})$, $v\in \bbL^2(\overline{\I_2})$ and $w\in \bbL^2(\overline{\I_3})$, the function $\bbE[f\mid\G_{i,j,k}]$ is the product of a function $\G_{i,j,k}\cap\overline{\I_1}$-measurable, a function $\G_{i,j,k}\cap\overline{\I_2}$-measurable and a function $\G_{i,j,k}\cap\overline{\I_3}$-measurable.

But we know that functions of the type $f=uvw$ generate a dense subalgebra of $\bbL^2(\A)$, so that we have proved that
\[
\G_{i,j,k}\subset \left(\G_{i,j,k}\cap\overline{\I_1}\right)\vee\left(\G_{i,j,k}\cap\overline{\I_2}\right)\vee\left(\G_{i,j,k}\cap\overline{\I_3}\right).
\]
which is the conclusion of the proposition since the other inclusion is trivial.
\end{proof}



\section{Limit law in the general case, for 2-dimensional fields}\label{dim2}
As shown in~\cite{MR3913270}, for a random field of martingale differences we have a CLT with convergence towards a mixture of normal laws (see Theorem below). In Section~\ref{product-type CLT}, for $f$ $\I$-measurable it was precised which mixtures can appear as limit laws (for definition of the factor $\I$ see Section~\ref{product-type factor}). Here we deal with the same question for the general case of a martingale difference $f\in \bbL^2$. We reduce our study to the case of (ergodic) $\Z^2$-actions. In many cases it is because we have not succeeded to extend the proofs to $d>2$.

As shown in Section~\ref{proj-mart}, if $f$ is a martingale difference, so is $\bbE[f \mid \I]$ and also $f -\linebreak \bbE[f \mid \I]$. Recall that the limit laws for the random field generated by $\bbE[f \mid \I]$ have been determined in Section~\ref{product-type CLT}. In Section~\ref{queue-property}, we give a sufficient condition guaranteeing convergence of the random field generated by $f - \bbE[f \mid \I]$ to a normal law. We show in Section~\ref{convolution} that, under the same condition the random field generated by $f$ is the convolution of the preceding ones.

In Section~\ref{entropie} we establish the result announced at the end of Section~\ref{proj-mart}.

Eventually in Section~\ref{nnormal_ex} we give an example of a field of martingale differences generated by $f - \bbE[f \mid \I]$ where the limit law is not normal. It remains an open question which mixtures of normal laws can appear as limits in the CLT for\linebreak $f - \bbE[f \mid \I]$.

In all this section, $f\circ T_{i,j}$ is a field of martingale differences adapted to a completely commuting filtration $\F_{i,j}$.

\subsection{Limit law for an increment orthogonal to the factor of product type}\label{queue-property}

Recall that $f\circ T_{i,j}$ is a field of martingale differences and as shown in Section~\ref{proj-mart}, $(f - \bbE[f \mid \I])\circ T_{i,j}$ are martingale differences as well.

Let us begin b recalling~\cite[Theorem~1]{MR3913270} which gives information on the limit law in the CLT. It will be stated and used here for $d=2$ but extends to any finite dimension.
\begin{theo}\label{thmdv}
When $\min\{m,n\}\to\infty$ the random variables $\frac1{\sqrt{mn}}\sum_{i=1}^m\sum_{j=1}^n \linebreak f\circ T_{i,j}$ converge in distribution to a law with characteristic function $\bbE[\exp(-\eta^2t^2/2)]$ where $\eta^2$ is a positive random variable such that $\bbE[\eta^2]=\|f\|_2^2$. The random variables $\frac1{mn}\sum_{i=1}^m\left(\sum_{j=1}^n f\circ T_{i,j}\right)^2$ converge in distribution to $\eta^2$.
\end{theo}

\begin{enonce*}[remark]{Comment on Theorem~\ref{thmdv}}
In fact, by the ergodic theorem,
\[
\lim_{m\to\infty}\frac1{m}\sum_{i=1}^m\frac1n\left(\sum_{j=1}^n f\circ
T_{i,j}\right)^2
\]
exists for each $n$ and the distribution of $\eta^2$ is the
limit in distribution of this quantity when $n\to\infty$. (This can be seen as well in the proof of Theorem~\ref{thmdv} or as a~consequence of it.)
\end{enonce*}
\begin{prop}\label{queue-condition}
Let $f\in \bbL^2\ominus \bbL^2(\I)$ be a martingale difference. If, moreover, $f\in\bbL^4$ and
\begin{equation}\label{eq:hyp_pour_avoir_loi_normale}
\lim_{\ell\to\infty}
\norm{
\E{f\mid\F_{\infty, -\ell}\vee \I_1 }}_2=0.
\end{equation}
then for $\min\{m,n\} \to\infty$, $(1/\sqrt{mn}) \sum_{i=1}^m\sum_{j=1}^n f\circ T_{i,j}$ converge in distribution to a centered normal law with variance $\bbE[f^2]$.
\end{prop}



\begin{rema}
We have $\I_2\subset \F_{\infty, -\ell}$ (for every $\ell$) hence $\I \subset
\F_{\infty, -\ell}\vee \I_1$. Therefore $\|\bbE[f\mid\I]\|_2 \leq \|\bbE[f\mid
\F_{\infty, -\ell}\vee \I_1]\|_2$ and, for any $f\in \bbL^2$, \eqref{eq:hyp_pour_avoir_loi_normale} implies $\bbE[f\mid\I] =0$.
\end{rema}
\begin{rema}
For any $f\in\bbL^2$, by the martingale convergence theorem, the sequence
\[
\left((\bbE[f\mid \Fca_{\infty,-\ell}\vee
\I_1])^2\right)
\]
converges in $\bbL^1$ hence is uniformly
integrable.
\end{rema}

As a consequence the property~\eqref{eq:hyp_pour_avoir_loi_normale} is equivalent to the convergence of the sequence $(\bbE[f\mid\F_{\infty, -\ell}\vee \I_1])_{\ell\geq 1}$ to zero in probability. Similarly, for $f\in \bbL^4$, condition~\eqref{eq:hyp_pour_avoir_loi_normale} implies that
\begin{equation}\label{eq:hyp_pour_avoir_loi_normale_4}
\lim_{\ell\to\infty}
\norm{
\E{f\mid\F_{\infty, -\ell}\vee \I_1 }}_4=0.
\end{equation}
\begin{rema}
Of course, condition~\eqref{eq:hyp_pour_avoir_loi_normale} implies
\begin{equation}\label{eq:condition_esperance_cond_avecF_infty_I1}
\E{f\mid \F_{\infty, -\infty}\vee \I_1}=0.
\end{equation}
In Section~\ref{trivial-queue} we give an example where
\[
\bigcap_\ell \F_{\infty, -\ell}\vee \I_1\neq\F_{\infty, -\infty}\vee \I_1,
\]
showing that~\eqref{eq:condition_esperance_cond_avecF_infty_I1} can be satisfied without~\eqref{eq:hyp_pour_avoir_loi_normale}. But we do not know if~\eqref{eq:condition_esperance_cond_avecF_infty_I1} is sufficient in order to obtain the conclusion of Proposition~\ref{queue-condition}.
\end{rema}
\begin{proof}[Proof of Proposition~\ref{queue-condition}]
Define
\begin{equation}\label{eq:def_Vmn}
V_{m,n}=\frac 1m\sum_{i=1}^m \left(\frac 1{\sqrt n}\sum_{j=1}^n f\circ T_{i,j}\right)^2.
\end{equation}
Following the theorem just recalled above, it is sufficient to prove that
\[
\lim_{n}\lim_{m} V_{m,n}=\norm{f}_2^2.
\]

By the ergodic theorem
\[
\lim_m V_{m,n}= \E{\left(\frac 1{\sqrt n}\sum_{j=1}^n f\circ T_{0,j}\right)^2\,\middle|\,\I_1}
\]
and the square terms give the expected limit; indeed, again by ergodic theorem,
\[
\lim_n \E{\frac 1{n}\sum_{j=1}^n f^2\circ T_{0,j}\,\middle|\,\I_1}=\E{\E{f^2\,\middle|\,\I_2}\,\middle|\,\I_1}
\]
and $\bbE[\bbE[f^2\mid\I_2]\mid\I_1]=\bbE[f^2]$ since the $\sigma$-algebras $\I_1$ and $\I_2$ are independent.

Thus it remains to prove that
\begin{equation}\label{eq:termes-croises}
\lim_n\E{\frac1n\sum_{1\leq j< k\leq n}f\circ T_{0,j}\;f\circ T_{0,k}\,\middle|\,\I_1} =0\quad\text{in norm }\norm{\cdot}_2.
\end{equation}
We'll need to work with order two moments of these sums, which are finite thanks to the assumption that $f\in\bbL^4$.

Let us write, for $0<\ell<n$,
\begin{multline*}
\frac1n\sum_{1\leq j< k\leq n}f\circ T_{0,j}\;f\circ T_{0,k} \\
\begin{aligned}
&=
\frac1n \sum_{k=\ell+1}^{n} f \circ T_{0,k}
\sum_{j=1}^{k-\ell} f \circ T_{0,j} +
\frac1n \sum_{k=2}^{n} f \circ T_{0,k}
\sum_{j=(k-\ell+1)\vee 1}^{k-1} f \circ T_{0,j}\\
&=: I + II.
\end{aligned}
\end{multline*}

The sequence $(f \circ T_{0,k} \sum_{j=(k-\ell+1)\vee 1}^{k-1} f\circ T_{0,j})_{k>1}$ is a sequence of square integrable martingale differences (remember that $f\in \bbL^4$), adapted to the filtration $(\F_{0,k})$. Pythagoras followed by Cauchy--Schwarz gives
\begin{align*}
\norm{II}_2^2&= \frac1{n^2} \sum_{k=2}^{n} \norm{ f \circ T_{0,k}
\sum_{j=(k-\ell+1)\vee 1}^{k-1} f \circ T_{0,j}}_2^2\\
&\leq \frac1{n^2} \sum_{k=2}^{n}
\norm{f\circ T_{0,k}}_4^2\norm{\sum_{j=(k-\ell+1)\vee 1}^{k-1} f \circ T_{0,j}}_4^2.
\end{align*}
Burkholder's inequality (see for example~\cite{MR2472010}) gives
\[
\norm{\sum_{j=(k-\ell+1)\vee 1}^{k-1} f \circ T_{0,j}}_4^2\leq3\sum_{j=(k-\ell+1)\vee 1}^{k-1}\norm{f\circ T_{0,j}}_4^2\leq 3\ell\norm{f}_4^2
\]
hence
\begin{equation}\label{eq:controle_norme_II}
\norm{\E{II\mid\I_1}}_2\leq\norm{II}_2\leq \sqrt{3}\norm{f}_4^2\sqrt{\frac{\ell}{n}}.
\end{equation}
The sequence $(f \circ T_{0,k}
\sum_{j=1}^{k-\ell} f \circ T_{0,j})_{k>\ell}$ is a sequence of square integrable martingale differences (remember that $f\in \bbL^4$), adapted to the filtration $(\F_{0,k})$. Applying Proposition~\ref{pro} to the $\Z$-action $(T_{k,0})$, $\I=\I_1$, $\F=\F_{\infty,k}$ and $\C$ trivial, we obtain that conditional expectations with respect to $\I_1$ and $\F_{\infty,k}$ commute. Moreover, since $\I_1\subset\F_{0,\infty}$, conditional expectations with respect to $\I_1$ and $\F_{0,\infty}$ commute. Using the commuting property of the filtration $(\F_{i,j})$, we affirm that conditional expectations with respect to $\I_1$ and $\F_{0,k}$ commute. As a consequence the sequence
\[
\left(\E{f \circ T_{0,k}
\sum_{j=1}^{k-\ell} f \circ T_{0,j}\,\middle|\,\I_1}\right)_{k>\ell}
\]
is a sequence of martingale differences, and in particular it is an orthogonal sequence in $ \bbL^2$. Using successively Pythagoras, invariance of the measure under $T_{0,k}$, properties of the conditional expectation, Cauchy--Schwarz and Burkholder, we write
\begin{align*}
\norm{\E{I\mid\I_1}}_2^2&= \frac1{n^2} \sum_{k=\ell+1}^{n}
\norm{\E{f \circ T_{0,k} \sum_{j=1}^{k-\ell} f \circ T_{0,j}
\,\middle|\,\I_1}}_2^2\\
&=\frac1{n^2} \sum_{k=\ell+1}^{n}
\norm{\E{f \sum_{j=1-k}^{-\ell} f \circ T_{0,j}
\,\middle|\,\I_1}}_2^2\\
&\leq\frac1{n^2} \sum_{k=\ell+1}^{n}
\norm{\E{f\mid\F_{\infty,-\ell}\vee\I_1} \sum_{j=1-k}^{-\ell} f \circ T_{0,j} }_2^2\\
&\leq\norm{\E{f \,\middle|\,\F_{\infty,-\ell}\vee\I_1}}_4^2 \ \frac1n\sum_{k=\ell+1}^n\norm{\frac1{\sqrt n}\sum_{j=1-k}^{-\ell} f \circ T_{0,j}}_4^2\\
&\leq 3\norm{\E{f \,\middle|\,\F_{\infty,-\ell}\vee\I_1}}_4^2\norm{f}_4^2.
\end{align*}

This estimation, associated with~\eqref{eq:controle_norme_II}, gives
\begin{multline*}
\norm{\E{\frac1n\sum_{1\leq j< k\leq n}f\circ T_{0,j}\;f\circ T_{0,k}\,\middle|\,\I_1}}_2
\\
\leq\sqrt3\left(\norm{\E{f \,\middle|\,\F_{\infty,-\ell}\vee\I_1}}_4\,\norm{f}_4+ \norm{f}_4^2\sqrt{\frac{\ell}{n}}\right),
\end{multline*}
hence for each $\ell$,
\begin{equation*}
\limsup_{n\to\infty}\norm{\E{\frac1n\sum_{1\leq j< k\leq n}f\circ T_{0,j}\;f\circ T_{0,k}\,\middle|\,\I_1}}_2\leq \sqrt{3}
\norm{\E{f \,\middle|\,\I_{1}\vee \F_{\infty, -\ell}}}_4\norm{f}_4,
\end{equation*}
and, thanks to~\eqref{eq:hyp_pour_avoir_loi_normale_4}, we conclude that~\eqref{eq:termes-croises} is true.
\end{proof}
\subsection{Entropy condition for convergence to a normal law}\label{entropie}

\begin{theo}\label{pos-entropy}
There exists a martingale difference $f\in \bbL^2$ with non normal limit in the CLT if and only if $T_{1,0}$ is of positive entropy in $\I_2$ and $T_{0,1}$ is of positive entropy in $\I_1$.
\end{theo}

\begin{proof}
1. Suppose that $T_{0,1}$ is of positive entropy in $\I_1$ and $T_{1,0}$ is of positive entropy in $\I_2$. For a measure preserving and bimeasurable transformation $S$, positive entropy implies existence of a nontrivial i.i.d.\ sequence of the form $h\circ S^i$. In the factor of product type given by $\I= \I_1\vee \I_2$ we are in the situation of Theorem~\ref{thm:TLC_product} and we thus get a non trivial field of Wang--Woodroofe type.
%Quel theorem 2? demander à l'auteur.

2. Suppose that the transformation $T_{0,1}$ is of zero entropy on $\I_1$. Then $\I_1$ is an invariant sub-$\sigma$-algebra of the Pinsker $\sigma$-algebra for $T_{0,1}$ hence by~\cite[Theorem~2]{MR928378} we have that for any integrable and $\F_{\infty,0}$-measurable function $g$, $\bbE[g\mid \F_{\infty, -\ell}\vee \I_1] = \bbE[g\mid \F_{\infty, -\ell}]$ for all $\ell >0$. Hence, if $g$ is a square integrable martingale difference, we have
\begin{equation*}
\E{g \,\middle|\,\F_{\infty, -\ell}\vee \I_1} = 0 \quad\text{for all } \ell >0.
\end{equation*}
Moreover, we know that $\I_2 \subset \F_{\infty,-1}$, so that $\I\subset \F_{\infty, -1}\vee \I_1$ and we have also $\bbE[g\mid\I]=0$. Thus, by Proposition~\ref{queue-condition}, we know that if $f$ is a martingale difference with finite fourth moment, then the CLT applies with a normal limit. For a square integrable martingale difference $f$, we write $f$ as a sum of a bounded martingale difference and a rest small in norm $\|\cdot\|_2$ (cf.~\cite[Theorem~4.2]{MR233396}); we conclude that the CLT applies to $f$ with a normal limit.
\end{proof}

Notice that Theorem~\ref{pos-entropy} improves the result in~\cite{MR3427925} which tells that the limit distribution is normal as soon as one of the $\sigma$-algebra $\I_1$ or $\I_2$ is trivial. However, the result in~\cite{MR3427925} applies to all $d>1$ while here it applies to $d=2$ only. The case of $d>2$ remains open.


\subsection{Remark on the decomposition given by the projection on the factor of product type}\label{convolution}

If $f$ is a martingale difference, we studied in Section~\ref{product-type CLT} the CLT for the martingale difference $\bbE[f \mid\I]$ and in Section~\ref{dim2} the CLT for the martingale difference $f-\bbE[f \mid\I]$. Using the theorem recalled p.~\pageref{thmdv}, it can be seen that if the limit law for $f-\bbE[f \mid\I]$ is normal, then the limit law for $f$ is the convolution of the limit laws for $\bbE[f \mid\I]$ and $f-\bbE[f \mid\I]$.

\subsection{Example of a non-normal limit law, in the orthocomplement of the product factor}\label{nnormal_ex}
Here we give an example of an ergodic $\Z^2$-action and a martingale difference $f$ which is orthogonal to the factor $\I$ but for which the limit distribution in the CLT is not normal.

Let us denote by $(\Omega,\B,\mu,S)$ the Bernoulli scheme $\bigl(\frac12,\frac12\bigr)$ on the alphabet $\{-1,1\}$. This means that $\Omega$ is the space of bilateral sequences of $-1$ or $1$, equipped with the product $\sigma$-algebra, the probability measure which makes the coordinate maps i.i.d.\ with law $\bigl(\frac12,\frac12\bigr)$, and the shift $S$.

Another dynamical system is $Z=\{-1,1\}$ equipped with the permutation $U$ and the uniform probability.

On the product probability space $\Omega\times\Omega\times Z$, consider the $\Z^2$-action $T$ defined by
\[
T_{i,j}(\omega,\omega',z)=\left(S^i\omega,S^j\omega',U^{i+j}z\right).
\]
This space is equipped with the \emph{natural} filtration $(\F_{i,j})$:
\[
\F_{i,j}=\sigma\left(X_k,Y_\ell,z ; k\leq i, \ell\leq j\right)=\sigma\left(X_k ; k\leq i\right)\otimes\sigma\left(Y_\ell ; \ell\leq j\right)\otimes\P(Z).
\]
where of course $X_k((\omega_n)_{n\in\Z})=\omega_k$ et $Y_\ell((\omega'_m)_{m\in\Z})=\omega'_\ell$ are the two independent Bernoulli processes.

The $\sigma$-algebra $\I_1$ of $T_{1,0}$-invariants is the $\sigma$-algebra of events depending only on the second coordinate:
\[
\I_1= (\text{trivial $\sigma$-algebra of $\Omega$})\otimes\B\otimes(\text{trivial $\sigma$-algebra of $Z$}).
\]
And symmetrically for the $\sigma$-algebra $\I_2$ of $T_{0,1}$-invariants. So that we have
\[
\I=\B\otimes\B\otimes(\text{trivial $\sigma$-algebra of $Z$}).
\]

We consider now the random variable $f(\omega,\omega',z)=X_0(\omega)Y_0(\omega')z$.

Since the expectation of $z$ is zero, we have $\bbE[f\mid\I]=0$. Associated to the $\Z^2$-action and the natural filtration, the function $f$ is a martingale difference.

But it is easy to calculate the limit distribution of $\bigl(\frac1{\sqrt{nm}}\sum_{i=1}^n\sum_{j=1}^m f\circ{T_{i,j}}\bigr)$. Indeed,
\[
\frac1{\sqrt{nm}}\sum_{i=1}^n\sum_{j=1}^m f\circ{T_{i,j}}(\omega,\omega',z)=\left(\frac1{\sqrt{n}}\sum_{i=1}^n(-1)^iX_i(\omega)\right)\left(\frac1{\sqrt{m}}\sum_{j=1}^m(-1)^jY_j(\omega')\right)z.
\]
The limit distribution is the distribution of the product of two independent $\clN(0,1)$ random variables, which is a non-normal Bessel law.

\subsection{Remark on the asymptotic condition insuring a normal law}\label{trivial-queue}

When looking to Proposition~\ref{queue-condition} the following question appears naturally. Is it true that
\begin{equation}\label{eq:bad_limit}
\bigcap_{\ell\geq0}\left(\F_{\infty,-\ell}\vee\I_1\right) = \F_{\infty,-\infty}\vee\I_1?
\end{equation}

We give here a negative answer by the construction of an example. (Note however that in this example the transformation $T_{2,0}$ is the identity so there will not exist any non zero martingale difference.)

The space is the bidimensional torus $\bbT^2=\R^2/\Z^2\simeq[0,1\mathclose{[}^2$ equipped with the Borel $\sigma$-algebra and the Lebesgue measure. The transformation $T_{0,1}$ is the automorphism defined by the matrix $
(\begin{smallmatrix} 
3&1\\2&1
\end{smallmatrix})$, that is $T_{0,1}(x,y)=(3x+y,2x+y)$, and the transformation $T_{1,0}$ is the translation $T_{1,0}(x,y)=(x+\frac12,y)$.

Transformations $T_{0,1}$ et $T_{1,0}$ are commuting:
\[
T_{1,0}T_{0,1}(x,y)=T_{0,1}T_{1,0}(x,y)=\left(3x+2y+\frac12,2x+y\right),
\]
so they generate a $\Z^2$-action denoted by $T$.

Let us denote by $\P$ the partition $\{[0,\frac12\mathclose{[}\times[0,1\mathclose{[},[\frac12,1\mathclose{[}\times[0,1\mathclose{[}\}$ of the torus, and by $(\F_j)$ the filtration generated by the transformation $T_{0,1}$ and this partition:
\[
\F_j=\sigma\left(T_{0,-k}\left(\P\right),\ k\leq j\right)\;.
\]
We know that the measure preserving dynamical system $(\bbT^2,T_{0,1})$ has the Kolmogorov property, thus the limit $\sigma$-algebra $\F_{-\infty}=\bigcap_{j\in\Z}\F_j$ is trivial (modulo the measure).

Note also that the partition $\P$ is invariant under the transformation $T_{1,0}$.

Define for all $i,j\in\Z$, $\Fca_{i,j}=\Fca_j$. Then $(\Fca_{i,j})_{i,j\in\Z}$ is a completely commuting invariant filtration. Moreover, $\Fca_{\infty,-\infty}\vee\I_1=\Fca_{-\infty}\vee\I_1=\I_1$.

Finally the following lemma shows that property~\eqref{eq:bad_limit} is not satisfied.

\begin{lemm}
For all $\ell\in\Z$, the $\sigma$-algebra $\F_{\infty,\ell}\vee\I_1$ is the whole Borel algebra.
\end{lemm}

\begin{proof}
Let us show that, for any $\ell$, the $\sigma$-algebra generated by $\I_1$ and the partition $T_{0,-\ell}(\P)$ is the whole Borel algebra.

For each integer $n$, consider
\[
\begin{pmatrix}3&1\\2&1
\end{pmatrix}^n=
\begin{pmatrix}a_n&b_n\\c_n&d_n
\end{pmatrix}
\]

A straightforward induction shows that, for all $n\in\Z$, the number $a_n$ is odd and the number $c_n$ is even.

Any measurable function of the two-dimensional variable $(2x,y)$ is $\I_1$ measurable and the map $(x,y)\mapsto\ind{[0,\frac12\mathclose{[}}(a_\ell x+b_\ell y)$ is $T_{0,-\ell }(\P)$-measurable. By multiplication, we obtain that the map $(x,y)\mapsto\ind{[0,\frac12\mathclose{[}}(x)$ is $T_{0,-\ell }(\P)\vee\I_1$-measurable. Using the representation $\bbT=[0,1[$ we write
\[
\exp(2i\pi x)=\exp\left(2i\pi\frac12(2x\modd 1)\right)\ind{\left[0,\frac12\right[}(x)-\exp\left(2i\pi\frac12(2x\modd 1)\right)\ind{\left[\frac12,1\right[}(x).
\]
This shows that any character of the two dimensional torus is $T_{0,-\ell }(\P)\vee\I_1$ measurable, proving that this $\sigma$-algebra is the whole Borel algebra.
\end{proof}
\begin{rema}
The idea behind the previous construction owes a great deal to an example attributed to Jean-Pierre Conze and communicated to us by Jean-Paul Thouvenot: if we denote by $\B(\alpha)$ the $\sigma$-algebra of Borel subsets of the one-dimensional torus invariant by the translation $x\mapsto x+\alpha$, we have that the algebra $
\bigcap_{n\in\N}\B(2^{-n})$ and $\bigcap_{n\in\N}\B(3^{-n}) $ are trivial (modulo the Lebesgue measure), but for each $n$ the $\sigma$-algebra $
\B(2^{-n})\vee\B(3^{-n}) $ is the whole Borel algebra.
\end{rema}




\begin{thebibliography}{99}
\bibitem{MR192027} Billingsley, Patrick Ergodic theory and information, Wiley Series in Probability and Mathematical Statistics, John Wiley \& Sons, 1965, xiii+195 pages | MR | Zbl
\bibitem{MR233396} Billingsley, Patrick Convergence of probability measures, John Wiley \& Sons, 1968 | MR | Zbl
\bibitem{MR3504508} Cuny, Christophe; Dedecker, Jérôme; Volný, Dalibor A functional CLT for fields of commuting transformations via martingale approximation, Zap. Nauchn. Semin. (POMI), Volume 441 (2015), pp. 239-262 | DOI | MR | Zbl
\bibitem{MR3869881} Giraudo, Davide Invariance principle via orthomartingale approximation, Stoch. Dyn., Volume 18 (2018) no. 6, 1850043, 29 pages | DOI | MR | Zbl
\bibitem{MR4754186} Giraudo, Davide Deviation inequality for Banach-valued orthomartingales, Stochastic Processes Appl., Volume 175 (2024), 104391, 20 pages | DOI | MR | Zbl
\bibitem{MR1958753} Glasner, Eli Ergodic theory via joinings, Mathematical Surveys and Monographs, 101, American Mathematical Society, 2003 | DOI | MR | Zbl
\bibitem{MR262464} Parry, William Entropy and generators in ergodic theory, Mathematics Lecture Note Series, W. A. Benjamin, Inc., 1969, xii+124 pages | Zbl | MR
\bibitem{MR2472010} Rio, Emmanuel Moment inequalities for sums of dependent random variables under projective conditions, J. Theor. Probab., Volume 22 (2009) no. 1, pp. 146-163 | DOI | MR | Zbl
\bibitem{MR928378} Volný, Dalibor Martingale decompositions of stationary processes, Yokohama Math. J., Volume 35 (1987) no. 1-2, pp. 113-121 | MR | Zbl
\bibitem{MR3427925} Volný, Dalibor A central limit theorem for fields of martingale differences, Comptes Rendus. Mathématique, Volume 353 (2015) no. 12, pp. 1159-1163 | DOI | MR | Zbl | Numdam
\bibitem{MR3913270} Volný, Dalibor On limit theorems for fields of martingale differences, Stochastic Processes Appl., Volume 129 (2019) no. 3, pp. 841-859 | DOI | MR | Zbl
\bibitem{MR3222815} Wang, Yizao; Woodroofe, Michael A new condition for the invariance principle for stationary random fields, Stat. Sin., Volume 23 (2013) no. 4, pp. 1673-1696 | DOI | MR | Zbl
\end{thebibliography}
\end{document}
